\section{问题三模型建立及求解}

% 后续问题模板：先说明真实继承，再说明问题三新增的对象、条件或数学困难。
% 不复制问题一/二的完整推导；若问题三独立，直接说明新的研究对象和数学缺口。

\subsection{模型建立与扩展}
% 建议顺序：继承的 current 关系 → 新增变量/状态 → 新目标或判据 → 新约束来源 → 最终可计算模型。
% 决定性推导、边界和约束不得因“前文已有模型”而省略；未变化部分用准确引用承接。
% 优化类仍应先说明决策对象和目标函数现实含义，并把目标函数放在约束大括号外。

\subsection{模型求解}
% 从问题三最终模型的计算结构或新增困难进入 solver。
% 写清继承/更换求解方法的原因、本题化编码、约束处理、初值/参数/精度/终止条件和输出映射。
% 算法步骤或伪代码只有在真实多阶段、循环、分支或筛选逻辑需要展示时才启用。

\subsection{求解结果}
% 给出问题三的 current 高精度核心答案、推荐方案或判定，并在核心图表邻近解释决定性特征。
% 结果段应直接回答问题三；不能只停在算法收敛、目标值或“结果如图所示”。

\subsection{结果的分析与验证}
% 仅在存在真实风险和 current 证据时启用。验证前先说明主结果已回答什么、仍受什么因素影响。
% 每项检验闭合：风险/主张 → 改变什么 → 保持什么 → 指标 → 变化范围 → 结论及有效边界。
% 若问题三简单且没有独立验证必要，删除整个小节。
